\section{Full Pairing Proof with Quotient Batching}
\label{sec:iop}

We now combine the polynomial equations from Section~\ref{sec:op-polys} into a complete two-round public-coin interactive proof, further reducing verification cost by batching all quotients into a single accumulated polynomial.

\subsection{Batching Polynomial Identity Testing}\label{sec:iop:batching}

All $n$ polynomial equations from Section~\ref{sec:op-polys} are batched into a single verification. The prover commits to remainders $\{R_i\}_{i=0}^{n-1}$ and receives challenge $z$. Using $z$, the prover forms the accumulated quotient $Q_\text{acc} = \sum_{i=0}^{n-1} z^i \cdot Q_i$. The verifier then samples a second challenge point $c$ and checks:

\begin{equation}
    \label{eq:iop-batched}
    \sum_{i=0}^{n - 1}z^i \cdot [A_i(c) \cdot B_i(c) \cdots X_i(c) - R_i(c)] = Q_\text{acc}(c) \cdot P_{12}(c)
\end{equation}

This reduces $n$ polynomial verifications (degrees 38--76, see Tables~\ref{tbl:op:ml-0bit}--\ref{tbl:op:finx}) to a single evaluation of a degree-76 polynomial.

\subsection{Prover's Algorithms}

The prover computes quotients $Q_i$ and remainders $R_i$ for each equation in Section~\ref{sec:op-polys}. For each step, it calls \textsc{PolyMulDiv}, which multiplies the input polynomials and performs standard Euclidean division by the irreducible $P_{12}$. After receiving challenge $z$, the prover forms $Q_\text{acc} = \sum_{i} z^i Q_i$ and sends $\mathcal{R} = \{R_i\}$ and $Q_\text{acc}$ to the verifier. The full prover procedure is given in Algorithm~\ref{alg:2r:prover}.

% \begin{algorithm}
% \caption{PolyMulDiv: Multiply polynomials and divide by $P_{12}$}
% \label{alg:polymuldiv}
% \begin{algorithmic}[1]
% \Require $P_1, \ldots, P_n \in \Fq[x]$; Divisor $P_{12} \in \Fq[x]$
% \Ensure $Q, R$ s.t.\ $\prod P_i = Q \cdot P_{12} + R$, $\deg(R) < \deg(P_{12})$
% \State $R \leftarrow \prod_{i=1}^{n} P_i$ \gComm{Polynomial product}
% \State Perform Euclidean division of $R$ by $P_{12}$ to get $(Q, R)$
% \State \Return $(Q, R)$
% \end{algorithmic}
% \end{algorithm}

\begin{algorithm}
\caption{Pairing Prover: Preparing polynomials}
\label{alg:2r:prover}
\begin{algorithmic}[1]
\Require $P \in \G_1, Q \in \G_2, c \in \Fe{12}, w \in \{0,1,2\}$
\Statex Internal: $W, W^2 \in \Fe{12}$ \Comment{$W$ is 27th root of unity}
\Ensure Arrays $\mathcal{Q}, \mathcal{R}$ of quotients and remainders
\Statex Write $s = 6t + 2 = \sum_{i=0}^{L-1} s_i 2^i$ where $s_i \in \{-1, 0, 1\}$ and $s_L = 1$
\State $c^{-1} \leftarrow 1 / c$ \Comment{Inverse Residue witness}
\State $f \leftarrow c^{-1}$ \Comment{Initialize with inverse residue witness}
\State $T \leftarrow Q$ \gComm{$T_i$ for each pair for multi-pairing}
\color{gray}
\Statex \small{Every operation involving $T, P \text{ or } Q$ is repeated for $T_i, P_i \text{ and } Q_i$ for multi-pairing.}
\color{black}

\Statex

\State \textbf{function} \texttt{PolyMulDiv} ({$R, P_1, \ldots, P_n \in \Fq[x]$}): \Return $Q, R \in \Fq[x]$ 
\State \qquad s.t.\ $\prod P_i = Q \cdot P_{12} + R$, $\deg(R) < \deg(P_{12})$

\Statex
\Statex \gText{Miller loop}
\For{$i = L - 2$ to $0$}
    \State $L_\text{dbl} \leftarrow l_{T,T}(P)$, $T \leftarrow [2]T$
    \If{$s_i = 0$}
		\State $Q, R \leftarrow {\textit{PolyMulDiv}}(f, f, L_\text{dbl})$ \Comment{add $L_{*i}$ for each $P_i, Q_i$}
    \ElsIf{$s_i = 1$}
        \State $L_\text{add} \leftarrow l_{T,Q}(P)$, $T \leftarrow T + Q$
        \State $Q, R \leftarrow {\textit{PolyMulDiv}}(f, f, c^{-1}, L_\text{dbl}, L_\text{add})$ \Comment{add $L_{*i}$ for each $P_i, Q_i$}
    \ElsIf{$s_i = -1$}
        \State $L_\text{add} \leftarrow l_{T,-Q}(P)$, $T \leftarrow T - Q$
        \State $Q, R \leftarrow {\textit{PolyMulDiv}}(f, f, c, L_\text{dbl}, L_\text{add})$ \Comment{add $L_{*i}$ for each $P_i, Q_i$}
    \EndIf
    \State $f \leftarrow R$, $\mathcal{Q}$.\textit{append}($Q$), $\mathcal{R}$.\textit{append}($R$)
\EndFor
\Statex


\State $Q_1 \leftarrow \pi_p(Q), Q_2 \leftarrow \pi_p^2(Q)$
\State $L_\pi \leftarrow l_{T,Q_1}(P)$
\State $T \leftarrow T + Q_1$
\State $L_{-\pi^2} \leftarrow l_{T,-Q_2}(P)$
\State $Q, R \leftarrow {\textit{PolyMulDiv}}(f, L_\pi, L_{-\pi^2})$ \Comment{add $L_{*i}$ for each $P_i, Q_i$}
\State $f \leftarrow R$, $\mathcal{Q}$.\textit{append}($Q$), $\mathcal{R}$.\textit{append}($R$)

\Statex

\Statex \gText{Final exponentiation: Cubic scaling + Frobenius mappings}

\If{$w = 0$}
    \State $Q, R \leftarrow {\textit{PolyMulDiv}}(f, c^{-q}, c^{q^2}, c^{-q^3})$
\ElsIf{$w = 1$}
    \State $Q, R \leftarrow {\textit{PolyMulDiv}}(f, W, c^{-q}, c^{q^2}, c^{-q^3})$
\ElsIf{$w = 2$}
    \State $Q, R \leftarrow {\textit{PolyMulDiv}}(f, W^2, c^{-q}, c^{q^2}, c^{-q^3})$
\EndIf
\State $\mathcal{Q}$.\textit{append}($Q$), $\mathcal{R}$.\textit{append}($R$) \gComm{Final $\mathcal{Q}$ and $\mathcal{R}$ accumulation, skip $f$}

\Statex

\State \Return $\mathcal{Q}, \mathcal{R}$
\end{algorithmic}
\end{algorithm}

\subsection{Verifier's Algorithms}

The verifier receives $P \in \G_1$, $Q \in \G_2$, residue witness $c$~\cite{OPP24}, the $W$ (27th root of unity), the remainder array $\mathcal{R}$, and $Q_\text{acc}$ from the prover. Using challenge $z$ from the first round and a freshly sampled $x$, it mirrors the prover's pairing loop while evaluating each polynomial at $x$, accumulating the RLC $e_\text{acc}$. The final check is Equation~\ref{eq:iop-batched}. The full verifier is given in Algorithm~\ref{alg:2r:verifier}.

% \begin{algorithm}
% \caption{EvalPoly: Evaluate polynomial product minus remainder via Horner's method}
% \label{alg:evalpolyprod}
% \begin{algorithmic}[1]
% \Require Polynomials $P_1, \ldots, P_n,\, R \in \Fq[x]$ of degree $\leq d$, evaluation point $z \in \Fq$
% \Ensure $e \in \Fq$ where $e = \prod_{i=1}^{n} P_i(z) - R(z)$

% \State \textbf{function} \texttt{EvalPoly} ({$R, P_1, \ldots, P_n \in \Fq[x]$}): \Return ($\prod_{i=1}^{n} P_i(z) - R(z) \in \Fq$)

% \State \textbf{function} \texttt{Eval} ({$P = (p_0, p_1, \ldots, p_d),\ z$}): \Return ($evaluation \in \Fq$)

% \State \gText{~~~~Uses Horner's method, body not included for brevity}

% \State $e \leftarrow \Call{Eval}{P_1, z}$ \gComm{Evaluate first polynomial}
% \For{$i = 2$ to $n$}
%     \State $e \leftarrow e \cdot \Call{Eval}{P_i, z}$ \gComm{Multiply in each subsequent evaluation}
% \EndFor
% \State $e \leftarrow e - \Call{Eval}{R, z}$ \gComm{Subtract remainder evaluation}
% \State \Return $e$
% \end{algorithmic}
% \end{algorithm}

\begin{algorithm}
\caption{Pairing Verifier: Evaluating remainders}
\label{alg:2r:verifier}
\begin{algorithmic}[1]
\Require $P \in \G_1, Q \in \G_2, c \in \Fe{12}, w \in \{0,1,2\}$, $\mathcal{R} \text{ array of remainders}$, $Q_\text{acc} \in \Fq[x]$
\Statex Internal: 27th root of unity, $W, W^2 \in \Fe{12}$, first round challenge, $z \xleftarrow{\$} \Fq$
\Ensure $pairing(P, Q) \stackrel{?}{=} 1$

\Statex Write $s = 6t + 2 = \sum_{i=0}^{L-1} s_i 2^i$ where $s_i \in \{-1, 0, 1\}$ and $s_L = 1$

\State $c^{-1} \leftarrow 1 / c$ \Comment{Inverse Residue witness}
\State $f \leftarrow c^{-1}$ \Comment{Initialize with inverse residue witness}
\State $T \leftarrow Q$ \gComm{$T_i$ for each pair for multi-pairing}
\State $x \xleftarrow{\$} \Fq$ \gComm{Sample evaluation point from field}
\State $e_\text{acc} \leftarrow 0$, $a_\text{rlc} \leftarrow 1$ \gComm{Evaluation accumulator and random linear combination multiplier}
\color{gray}
\Statex \small{Every operation involving $T, P \text{ or } Q$ is repeated for $T_i, P_i \text{ and } Q_i$ for multi-pairing.}
\color{black}

\Statex

\State \textbf{function} \texttt{EvalPoly} ({$R, P_1, \ldots, P_n \in \Fq[x]$}): \Return ($\prod_{i=1}^{n} P_i(z) - R(z) \in \Fq$)

\Statex
\Statex \gText{Miller loop}
\For{$i = L - 2$ to $0$}
    \State $L_\text{dbl} \leftarrow l_{T,T}(P)$, $T \leftarrow [2]T$
    \State $R \leftarrow \mathcal{R}.\textit{pop\_front}()$ \gComm{Fetch R from prover}
    \If{$s_i = 0$}
		\State $e_\text{acc} \leftarrow e_\text{acc} + a_\text{rlc} * \texttt{EvalPoly}(R, x, f, f, L_{\text{dbl}})$ \Comment{add $L_{*i}$ for each $P_i, Q_i$}
    \ElsIf{$s_i = 1$}
        \State $L_\text{add} \leftarrow l_{T,Q}(P)$, $T \leftarrow T + Q$
        \State $e_\text{acc} \leftarrow e_\text{acc} + a_\text{rlc} * \texttt{EvalPoly}(R, x, f, f, c^{-1}, L_\text{dbl}, L_\text{add})$ \Comment{add $L_{*i}$ for each $P_i, Q_i$}
    \ElsIf{$s_i = -1$}
        \State $L_\text{add} \leftarrow l_{T,-Q}(P)$, $T \leftarrow T - Q$
        \State $e_\text{acc} \leftarrow e_\text{acc} + a_\text{rlc} * \texttt{EvalPoly}(R, x, f, f, c, L_\text{dbl}, L_\text{add})$ \Comment{add $L_{*i}$ for each $P_i, Q_i$}
    \EndIf
    \State $f \leftarrow R$, $a_\text{rlc} \leftarrow a_\text{rlc} \cdot z$ \gComm{Update random linear combination factor}
\EndFor
\Statex

\State $R \leftarrow \mathcal{R}.\textit{pop\_front}()$ \gComm{Fetch R from prover}
\State $Q_1 \leftarrow \pi_p(Q), Q_2 \leftarrow \pi_p^2(Q)$
\State $L_\pi \leftarrow l_{T,Q_1}(P)$
\State $T \leftarrow T + Q_1$
\State $L_{-\pi^2} \leftarrow l_{T,-Q_2}(P)$
\State $e_\text{acc} \leftarrow e_\text{acc} + a_\text{rlc} * \texttt{EvalPoly}(R, x, f, L_\pi, L_{-\pi^2})$ \Comment{add $L_{*i}$ for each $P_i, Q_i$}
\State $f \leftarrow R$, $a_\text{rlc} \leftarrow a_\text{rlc} \cdot z$ \gComm{Update random linear combination factor}

\Statex

\Statex \gText{Final exponentiation: Cubic scaling + Frobenius mappings}
\State $R \leftarrow \mathcal{R}.\textit{pop\_front}()$ \gComm{Fetch R from prover}
\If{$w = 0$}
    \State $e_\text{acc} \leftarrow e_\text{acc} + a_\text{rlc} * \texttt{EvalPoly}(R, x, f, c^{-q}, c^{q^2}, c^{-q^3})$
\ElsIf{$w = 1$}
    \State $e_\text{acc} \leftarrow e_\text{acc} + a_\text{rlc} * \texttt{EvalPoly}(R, x, f, W, c^{-q}, c^{q^2}, c^{-q^3})$
\ElsIf{$w = 2$}
    \State $e_\text{acc} \leftarrow e_\text{acc} + a_\text{rlc} * \texttt{EvalPoly}(R, x, f, W^2, c^{-q}, c^{q^2}, c^{-q^3})$
\EndIf

\Statex

\State \Return $Q_\text{acc}(x) \stackrel{?}{=} e_\text{acc} \land R \stackrel{?}{=} 1$ \gComm{Polynomial identity and expected pairing output check}
\end{algorithmic}
\end{algorithm}

\subsection{Interactive Proof} \label{sec:iop:interactive}

\IOP{Two-Round Pairing IOP}{Verifier}{Prover}{
    \RDoes{Computes pairings with Algorithm~\ref{alg:2r:prover} and prepares\\ remainders and quotients arrays $\mathcal{R}$ and $\mathcal{Q}$}{1}
    \RtoL{Round I: Commits \textit{pairing inputs} and remainders array $\mathcal{R}$}{2.5}
    \LtoR{Challenge $z \xleftarrow{\$} \Fq$}{3.5}
    \RDoes{Computes $Q_{\text{acc}} = \sum_{i=0}^{n-1} z^i \cdot Q_i$ using $\mathcal{Q}$ and $z$}{4.5}
    \RtoL{Round II: Commit $Q_\text{acc}$}{5.5}
    \LDoes{Samples $x \xleftarrow{\$} \Fq$ and runs Algorithm~\ref{alg:2r:verifier} to test\\polynomial identities at $x$, Eq~\ref{eq:iop-batched}}{6.5}
    \LtoR{Accept/Reject}{8}
}

\subsection{Fiat-Shamir Transformation} \label{sec:iop:fs}

Applying the Fiat-Shamir transformation~\cite{fiatshamir} with random oracle $H$ to the two-round protocol of Section~\ref{sec:iop:interactive} gives a non-interactive proof with knowledge soundness $\delta_s + 3(m^2+1)2^{-\lambda}$~\cite{IOP-NI}, where $m=2$ is the number of oracle queries (one per round) and $\lambda$ is the challenge bit-length; Section~\ref{sec:iop:soundness} instantiates this bound.

\subsubsection{Non-Interactive Proof} \label{sec:iop:noninteractive}

The prover runs Algorithm~\ref{alg:2r:prover} to obtain $(\mathcal{Q}, \mathcal{R})$, derives the challenge
$$z = H(\text{\textit{pairing inputs}},\, c,\, W^w,\, \mathcal{R})$$
from the transcript instead of receiving it, computes $Q_{\text{acc}} = \sum_{i=0}^{n-1} z^i \cdot Q_i$, and outputs the proof $(\text{\textit{pairing inputs}},\, c,\, W^w,\, \mathcal{R},\, Q_\text{acc})$. The verifier recomputes $z$, derives the evaluation point $x = H(z,\, Q_\text{acc})$, and runs Algorithm~\ref{alg:2r:verifier}, testing the polynomial identities via Equation~\eqref{eq:iop-batched}.

The soundness bound of Theorem~\ref{thm:soundness} carries over because every prover-chosen value entering Equation~\eqref{eq:iop-batched} is absorbed before the challenge it must not depend on: $\mathcal{R}$ and the hinted witnesses $c, W^w$ before $z$, and $Q_\text{acc}$ with $z$ chained into the second query before $x$.

\subsubsection{Instantiating the Random Oracle}\label{sec:iop:fs:ro}

The protocol is agnostic to how $H$ is realised; what varies across execution environments is the cost of a \emph{commitment}, i.e.\ of absorbing one $\Fq$ element into the transcript. We therefore qualify the commitment counts of Section~\ref{sec:ops:comp-analysis} and Table~\ref{tab:garaga_performance} for our two instantiations.

\paragraph{CairoVM (Poseidon).} \label{sec:iop:fs:cairo}
Prover-supplied hints are passed as transaction calldata and absorbed into a running Poseidon~\cite{poseidon} sponge; $\Fq$ elements are represented by 96-bit limbs, with two limbs packed into each \texttt{Felt252} absorption. Both hashing cost and calldata therefore scale linearly with the number of committed elements, which is why the COMMITS column of Table~\ref{tab:garaga_performance} is reported alongside modular operations: halving the committed elements roughly halves the transcript-related share of VM steps and the hint calldata.

\paragraph{gnark (\texttt{api.Commit}).} \label{sec:iop:fs:gnark}
The two rounds map to the \textbf{two} sequential \texttt{api.Commit} calls in \texttt{performDeferredRingChecks} (Section~\ref{sec:polyring:mapping}): the first commits the limbs of the queued remainders and hinted witnesses, yielding $z$; the second commits $z$ together with the limbs of $Q_\text{acc}$, yielding $x$. The cost of \texttt{Commit} is backend-specific. In Groth16, gnark implements the commitment scheme of Belling et al.~\cite{BellingCCS23,gnarkcommit}: the committed wires are aggregated into a Pedersen-style $\G_1$ commitment, and the challenge is a hash-to-field of the commitment point, computed \emph{natively} by prover and verifier---no in-circuit hashing. Each \texttt{Commit} costs the prover two $\G_1$ MSMs of size equal to the number of committed wires (the commitment and its proof of knowledge); the proof grows by one $\G_1$ point per commitment plus a single proof of knowledge folded across commitments (64 bytes each, uncompressed); and verification gains one two-pairing product check, one hash-to-field, and one extra scalar multiplication and point addition in the public-input MSM.
In PLONK, committed wires reuse the existing wire-commitment machinery, so the marginal overhead per commitment is lower still, but each committed wire adds a constraint.

\paragraph{Impact.} The number of elements committed to determines hashing and calldata in CairoVM and the commitment-MSM size in gnark. Reducing the absolute number of committed elements, by 53\% through our consolidation (Section~\ref{sec:ops:comp-analysis}), directly reduces end-to-end cost regardless of the instantiation, without adding commitment rounds.


\subsection{Security Analysis}

Throughout this section, $z, x \in \mathbb{F}_q$ denote field element challenges (scalars),
while polynomials are elements of $\mathbb{F}_q[X]$ with $X$ as the formal indeterminate.
In particular, $z^j$ is a scalar coefficient, not a polynomial term.

\subsubsection{Algebraic Intuition}\label{sec:iop:alg-intuition}

Before the formal soundness proof it is instructive to consider why an incorrect
remainder is difficult to hide.  Suppose a malicious prover corrupts exactly one remainder:
$R_j' = R_j + \Delta_j$ with $\Delta_j \in \mathbb{F}_q[X]$, $\Delta_j \neq 0$,
$\deg(\Delta_j) \leq 11$.  Since $z$ is obtained \emph{after} the commitment to
$\mathcal{R}$, the prover cannot adjust any other $R_k$ to compensate for $\Delta_j$.
To successfully deceive the verifier, the prover must find a valid forged quotient
$Q'_\text{acc} \in \mathbb{F}_q[X]$ satisfying

$$\sum_{i=0}^{n-1} z^i(A_i - R_i') = Q'_\text{acc} \cdot P_{12}$$

All $R_i' = R_i$ except for $R_j' = R_j + \Delta_j$. Substituting $R_j' = R_j + \Delta_j$ gives
\begin{align*}
Q'_\text{acc} \cdot P_{12} &= Q_\text{acc} \cdot P_{12} - z^j\Delta_j(x) \\
Q'_\text{acc} &= Q_\text{acc} - \frac{z^j\Delta_j(x)}{P_{12}(x)}.
\end{align*}
Here $z^j \in \mathbb{F}_q$ is a nonzero scalar (with overwhelming probability over the
choice of $z$) and $\Delta_j(x) \in \mathbb{F}_q[X]$ is a nonzero polynomial of degree
at most~$11$.  Since $P_{12}$ is irreducible of degree~$12$ and
$\deg(z^j\Delta_j) \leq 11 < 12 = \deg(P_{12})$, the polynomial $P_{12}$ does not
divide $z^j\Delta_j$ in $\mathbb{F}_q[X]$; hence $z^j\Delta_j/P_{12}$ is not a
polynomial. Consequently no valid $Q'_\text{acc} \in \mathbb{F}_q[X]$ can satisfy the
equation for this single-corruption case.


The same degree argument extends immediately to an adversary who corrupts
\emph{multiple} remainders: for any scalar $z$, the weighted sum
$\sum_i z^i \Delta_i(x)$ is still a polynomial in $x$ of degree at most~$11$
(since each $z^i$ is a scalar and each $\deg(\Delta_i) \leq 11$), so it cannot
be divisible by $P_{12}$ of degree~$12$ unless it is the zero polynomial, which
requires every $\Delta_i = 0$.  The formal soundness proof below re-derives this
uniformly via Schwartz-Zippel.


\subsubsection{Soundness}\label{sec:iop:soundness}

\begin{theorem}[Soundness]\label{thm:soundness}
Let $n$ be the number of polynomial equations batched via a random linear combination into the bivariate polynomial identity:
$$P_{\mathrm{final}}(z, x) = \sum_{i=0}^{n-1} z^i \cdot \Bigl[\prod_{j} P_{i,j}(x) - R_i(x)\Bigr] - Q_{\mathrm{acc}}(x) \cdot P_{12}(x) = 0$$
Let the maximum degree in $x$ of the polynomial products be $d_x < 2^8$. For a prover committing to incorrect remainders, the probability of successful forgery $\delta_s$ over randomly chosen $(z,x) \in \mathbb{F}_q^2$ is bounded by:
$$\delta_s \leq \frac{n - 1 + d_x}{|\mathbb{F}_q|}$$
\end{theorem}

\begin{proof}
If the prover commits to an incorrect remainder $R_i(x)$, the evaluation yields a non-zero error polynomial $P_{\mathrm{err}}(z, x)$. The structure of $P_{\mathrm{final}}(z, x)$ establishes the following degree bounds for $P_{\mathrm{err}}(z, x)$:
\begin{itemize}
    \item \textbf{Degree in $z$:} at most $n - 1$ (due to the $z^i$ weighting).
    \item \textbf{Degree in $x$:} at most $d_x < 2^8$ (due to the polynomial products).
\end{itemize}
By the Schwartz-Zippel lemma, the probability that the non-zero polynomial $P_{\mathrm{err}}$ evaluates to zero is bounded by the total degree divided by the field size:
$$\Pr[P_{\mathrm{err}}(z, x) = 0] \leq \frac{\deg_z(P_{\mathrm{err}}) + \deg_x(P_{\mathrm{err}})}{|\mathbb{F}_q|} \leq \frac{n - 1 + d_x}{|\mathbb{F}_q|}$$
\end{proof}

\textbf{Concrete parameters for BN254}: With $n = 67$ equations and $d_x \leq 256$:
$$\delta_s \leq \frac{66 + 256}{2^{254}} < \frac{2^9}{2^{254}} = 2^{-245}.$$

\textbf{Non-interactive proof via Fiat-Shamir} \label{sec:iop:soundness:fs}: From Section~\ref{sec:iop:fs}, under the random oracle model the
knowledge soundness error becomes~\cite{IOP-NI}:
$$\delta_s' = \delta_s + 3(m^2 + 1) \cdot 2^{-\lambda}$$
where $m = 2$ (one query per round) and $\lambda = 254$ (challenge bit-length), giving
$$\delta_s' \leq 2^{-245} + 15 \cdot 2^{-254} < 2^{-244}.$$

\begin{remark}[Correctness of polynomial evaluation]\label{rem:overflow}
Theorem~\ref{thm:soundness} assumes that the verifier evaluates all polynomials at the
challenge point $x$ without arithmetic error.  The circuit implementation uses the
deferred-reduction inner-product optimisation of Section~\ref{sec:polyring:performance}
(Section~\ref{sec:emulated:flat}), which accumulates limb-wise products over $\mathbb{Z}$
before a single final reduction.  Correctness of this optimisation requires that no
intermediate limb value reaches the native field characteristic $p$.  We verify this
bound concretely for BN254 in Lemma~\ref{lem:overflow} (Section~\ref{sec:emulated:flat}).
\end{remark}

\subsubsection{Completeness}

For an honest prover executing Algorithm~\ref{alg:2r:prover} with valid inputs:

\begin{enumerate}
\item \textbf{Polynomial construction}: At each operation, the prover computes following polynomials:
   $$\prod_{j} P_{i,j} = Q_i \cdot P_{12} + R_i$$

\item \textbf{RLC formation}: Given challenge $z$, the prover forms:
   $$Q_{\text{acc}} = \sum_{i=0}^{n-1} z^i \cdot Q_i$$

\item \textbf{Verifier's check}: The verifier evaluates at point $x$:
   \begin{align*}
   \text{LHS} &= \sum_{i=0}^{n-1} z^i \cdot \left[\prod_{j} P_{i,j}(x) - R_i(x)\right]\\
   &= \sum_{i=0}^{n-1} z^i \cdot Q_i(x) \cdot P_{12}(x) \quad \text{(by construction)}\\
   &= P_{12}(x) \cdot \sum_{i=0}^{n-1} z^i \cdot Q_i(x)\\
   &= P_{12}(x) \cdot Q_{\text{acc}}(x) = \text{RHS}
   \end{align*}
\end{enumerate}

Since the equality holds for any evaluation point $x$, the protocol satisfies perfect completeness.

% https://github.com/shramee/cairo_pairing/blob/main/packages/bn254_u256/src/pairing/schzip/steps.cairo
